\section{P4: Programming Protocol-Independent Packet Processors}

\textit{Pat Bosshart et al., ACM SIGCOMM Computer Communication Review, 2014. This chapter follows the annotated paper and concentrates on the highlighted concepts and the technical examples used to explain them.}

\subsection{Motivation: Why P4?}

Software-Defined Networking (SDN) gives operators programmatic control over a network. Its key architectural separation is between the \textbf{control plane}, which decides how traffic should be handled, and the \textbf{forwarding/data plane}, which processes packets. A single control plane may manage multiple forwarding devices. A common, open, vendor-agnostic interface such as OpenFlow lets the controller manage devices from different vendors.

The limitation of early OpenFlow is that it explicitly names the packet-header fields and actions that switches support. OpenFlow 1.0 began with 12 matchable header fields; later versions expanded the set to 15, 36, 40, and 41 fields. The specification therefore became more complex while still not giving operators a general way to introduce new header formats. New encapsulations such as NVGRE, VXLAN, and STT illustrate why a fixed list of supported protocols is limiting.

The central idea of P4 is to let a programmer describe \emph{how a switch should parse and process packets}, rather than forcing the programmer to accept a fixed parser and fixed-function pipeline. P4 is a high-level language for programming protocol-independent packet processors and works alongside SDN control protocols such as OpenFlow. OpenFlow-like APIs can populate forwarding-table entries; P4 describes the packet-processing structure those tables belong to.

\subsection{The three goals of P4}

\begin{enumerate}
  \item \textbf{Reconfigurability in the field.} After a switch has been deployed, its packet parsing and processing should be redefinable. For example, the switch should be able to recognize a newly introduced header without requiring the design to have anticipated that header in advance.
  \item \textbf{Protocol independence.} The switch should not be tied to particular packet formats. The controller/programmer should be able to specify a parser that extracts named, typed fields and a collection of typed match--action tables that operate on those fields.
  \item \textbf{Target independence.} The programmer should describe the desired packet-processing behavior without depending on the underlying switch hardware. A compiler considers the target's capabilities when translating the P4 program into a target-specific configuration. The analogy in the paper is that a C programmer need not program directly for the internal details of a particular CPU.
\end{enumerate}

These goals require a balance between expressiveness and ease of implementation across different hardware and software switches. P4 is proposed as a ``sweet spot'': more flexible than fixed-function OpenFlow 1.0, but more constrained and easier to map to hardware than a fully general packet-processing environment such as Click.

\subsection{OpenFlow's limitation and the need for a higher-level language}

OpenFlow provides a common interface for populating forwarding tables, but its fixed parser and predefined match fields constrain which protocols the switch can recognize. Adding a new protocol repeatedly to the OpenFlow specification is not a scalable long-term approach. P4 instead allows the controller/programmer to describe the parser, the fields used by tables, and the processing/control flow between tables.

The paper's comparison is important:
\begin{itemize}
  \item \textbf{OpenFlow 1.x:} a common interface, but a predetermined set of header fields and a small set of predefined actions.
  \item \textbf{Click:} highly expressive and suitable for composing software packet-processing modules, but too unconstrained for easy inference of dependencies and mapping to dedicated hardware; it also does not directly provide the controller--switch table-population model described here.
  \item \textbf{P4:} an imperative description of packet-processing logic designed to be translated into a table-dependency representation and then mapped to different switch targets.
\end{itemize}

\subsection{Abstract forwarding model}

The abstract model processes an arriving packet through a \textbf{programmable parser} followed by multiple \textbf{match+action stages}. The stages may be arranged in series, in parallel, or as a combination of both. This generalizes the fixed sequence of stages assumed by the original OpenFlow model.

The model distinguishes two operations:
\begin{itemize}
  \item \textbf{Configure:} program the parser, define the order/arrangement of match+action stages, and specify which header fields each stage processes. Configuration determines which protocols the switch supports and how it can process packets.
  \item \textbf{Populate:} add or remove entries in the match+action tables already defined by configuration. Population determines the current policy applied to packets.
\end{itemize}

The paper treats configuration and population as separate phases and assumes that packet processing need not occur during configuration. However, implementations may support partial or full reconfiguration while continuing to forward packets, enabling upgrades without downtime. This configuration phase has little practical meaning for a fixed-function ASIC; in that case the compiler's task is mainly to check whether the chip can support the P4 program.

\subsubsection{Parser}

The parser examines packet headers, extracts their fields, and recognizes the sequence of headers. It therefore defines which protocols the switch can parse. The packet body is assumed to be buffered separately and is not available for matching in this model. The model does not assume any particular meaning for a header: once parsed, its fields are simply values on which match and action operations can work.

\subsubsection{Ingress and egress match+action processing}

After parsing, extracted fields are passed to match+action tables divided into ingress and egress processing.
\begin{itemize}
  \item \textbf{Ingress processing} can modify packet headers, select the egress port or ports, and select the output queue. Based on its result, a packet can be forwarded, replicated (for multicast, mirroring/span, or delivery to the control plane), dropped, or used to trigger flow control.
  \item \textbf{Egress processing} can perform per-instance modifications to the packet header, such as changes that differ across multicast copies.
\end{itemize}

Action tables such as counters and policers can be associated with a flow to track state across packets. Queueing is modeled as in OpenFlow: an action maps a packet to a queue, while the queue's service discipline (for example, minimum rate or deficit round robin (DRR)) is chosen as part of switch configuration.

\subsubsection{Metadata}

Packets can carry \textbf{metadata} between processing stages. Metadata is treated like a packet-header field for matching and actions, but it need not change the parsed representation of the packet. Examples include:
\begin{itemize}
  \item ingress port;
  \item transmit destination and queue;
  \item a timestamp used for packet scheduling;
  \item information passed from one table to another, such as a virtual-network identifier.
\end{itemize}

The paper also notes that action primitives could be extended for congestion-control protocols. For example, a switch could set the ECN bit under new conditions or implement a proprietary congestion-control mechanism with match+action tables; the detailed design of such primitives is outside the paper's scope.

\subsection{Table dependencies and parallel execution}

A packet-processing program must express dependencies between header fields and tables. A later table may depend on a field produced or modified by an earlier table, in which case those tables cannot safely execute in parallel. If two tables do not depend on each other's results, they may be eligible for parallel execution.

For example, the paper states that an IP routing table and an ARP table require sequential execution because of a data dependency between them. These dependencies can be represented by a \textbf{Table Dependency Graph (TDG)}:
\begin{itemize}
  \item each TDG node corresponds to a match+action table;
  \item edges represent dependencies induced by field inputs, actions, and control flow;
  \item dependency analysis helps determine where tables can be placed in the target pipeline and which tables may run in parallel.
\end{itemize}

Most programmers naturally describe algorithms using imperative constructs rather than dependency graphs. P4 therefore lets programmers express control flow imperatively, then uses a compiler to derive the TDG for dependency analysis and hardware mapping.

\subsection{Two-step compilation model}

The proposed compilation process has two conceptual steps:
\begin{enumerate}
  \item Translate the imperative P4 control program into an intermediate Table Dependency Graph (TDG), then analyze dependencies and opportunities for concurrency.
  \item Map the TDG to a particular switch target, allocating its available resources and generating target-specific configuration.
\end{enumerate}

The compiler thus separates the programmer's target-independent description from hardware-specific implementation decisions. The target may be a software switch, a multicore software switch, a network processor, a fixed-function switch ASIC, or a reconfigurable match-table pipeline.

\section{P4 language by example: the mTag design}

\subsection{Problem setting}

The motivating example is a Layer 2 network with top-of-rack (ToR) switches at the edge and a two-tier core. As the number of end hosts grows, the core's Layer 2 forwarding tables may overflow. MPLS could simplify core forwarding, but implementing a label-distribution protocol with multiple tags is considered daunting. PortLand is another option, but it requires rewriting MAC addresses—which may break existing network-debugging tools—and adding agents to respond to ARP requests.

The paper proposes a small illustrative scheme called \textbf{mTag}, combining hierarchical routing similar to PortLand with simple MPLS-like tags. The example is intentionally simplified to explain P4, rather than to specify a complete production switch.

\subsection{mTag representation and forwarding idea}

The mTag is a 32-bit label consisting of four one-byte fields:
\[
\texttt{up1},\;\texttt{up2},\;\texttt{down1},\;\texttt{down2}.
\]
It can encode a source route or a destination locator (similar to PortLand's pseudo-MAC). Each core switch only needs to inspect one byte of the tag, chosen according to the switch's position in the hierarchy and the direction of travel (up or down). In the example, the first ToR switch adds the tag, although the end-host network interface card could add it instead.

The purpose is to reduce the forwarding work required in the core: instead of looking up a full Layer 2 destination in a potentially large table, a core switch can forward using the relevant byte of the hierarchical tag.

\subsection{The four main P4 program components}

A P4 program describes four central components, together with the control program that connects their processing:
\begin{itemize}
  \item \textbf{Headers:} ordered fields, their widths, and constraints on field values.
  \item \textbf{Parsers:} the valid headers and valid sequences of headers in a packet.
  \item \textbf{Tables:} match+action tables, including the fields that may be matched and the actions that may be executed.
  \item \textbf{Actions:} functions composed from protocol-independent primitive operations.
  \item \textbf{Control program:} imperative logic that specifies the order in which tables are applied and the control flow between them.
\end{itemize}

\subsection{Header format declarations}

A header is defined by an ordered list of field names and their bit widths. Optional annotations can impose constraints on values or on the maximum lengths of variable-sized fields. The paper gives Ethernet and VLAN as examples:
\begin{itemize}
  \item Ethernet: destination address (48 bits), source address (48 bits), and EtherType (16 bits).
  \item VLAN: PCP (3 bits), CFI (1 bit), VLAN ID (12 bits), and EtherType (16 bits).
\end{itemize}

The mTag header can be added without changing the existing Ethernet/VLAN declarations. It contains four 8-bit fields (\texttt{up1}, \texttt{up2}, \texttt{down1}, \texttt{down2}) and a 16-bit EtherType. The field names communicate the two aggregation levels and the up/down forwarding directions.

\subsection{Programmable packet parser}

P4 assumes the target can implement a state machine that traverses packet headers, extracts field values, and selects the next state. The parser is expressed as transitions from one header to the next; a transition may be selected based on a field value in the current header.

In the mTag example, parsing begins at the start state and extracts Ethernet. The Ethernet EtherType determines whether the next header is VLAN, IPv4, or another supported case. The VLAN EtherType can indicate mTag (\texttt{0xaaaa}) or IPv4 (\texttt{0x800}); after mTag, the EtherType can again indicate IPv4. Parsing ends at an explicit stop state or when an unhandled case is encountered (which may be treated as an error).

The extracted headers are then passed to the match+action pipeline. The example parser has only four states, whereas real-world parsers can require many more; the paper notes that one previously described parser expands to over one hundred states.

\subsection{Match+action table specification}

The programmer specifies:
\begin{itemize}
  \item which header fields a table reads;
  \item the matching type for each field (for example, exact, range, or wildcard/ternary matching);
  \item the actions that may be selected when an entry matches;
  \item the maximum number of entries the table must support.
\end{itemize}

In the example, the edge switch matches the Layer 2 destination address and VLAN ID, then selects an mTag to insert. The table's \texttt{reads} attribute declares the match fields and their match types; its \texttt{actions} attribute declares the permitted actions; and \texttt{max\_size} states the required capacity. This information helps the compiler determine the amount and type of memory needed, such as TCAM or SRAM.

The paper's example mTag table has a maximum size of 20,000 entries. Other tables are used to check the source port, perform local switching, and validate the egress state. The source-check table verifies that an mTag is seen only on appropriate ports connected to core switches; an inappropriate mTag can be sent to the CPU, while a valid one is stripped and recorded in metadata. The local-switching table checks whether the destination is a locally connected host; if it misses, processing proceeds to the mTag table. The egress-check table checks that egress resolution is valid and prevents inappropriate retagging.

\subsection{Primitive and compound actions}

P4 provides primitive actions from which more complex action functions are constructed. The paper expects primitives within an action function to execute in parallel; a switch that cannot do so may emulate the same semantics.

The example \texttt{add\_mTag} action performs the following operations:
\begin{enumerate}
  \item Add an mTag header instance.
  \item Copy the VLAN header's EtherType into the new mTag's EtherType so that the encapsulated protocol is preserved.
  \item Set the VLAN EtherType to \texttt{0xaaaa}, marking that an mTag follows.
  \item Write the runtime parameters into \texttt{up1}, \texttt{up2}, \texttt{down1}, and \texttt{down2}.
  \item Set the metadata egress specification to the selected output port.
\end{enumerate}

The switch therefore inserts mTag after the VLAN header. The VLAN EtherType identifies the mTag, and the mTag EtherType records the protocol that follows it. The paper notes that the inverse operation—removing an mTag—and the table that invokes it on edge switches are not shown in full.

The primitive actions listed are:
\begin{itemize}
  \item \textbf{set field:} assign a value to a header field; masked sets are supported;
  \item \textbf{copy field:} copy one field into another;
  \item \textbf{add header:} make a header instance and its fields valid;
  \item \textbf{remove header:} pop a header and its fields from the packet;
  \item \textbf{increment:} increment or decrement a field;
  \item \textbf{checksum:} compute a checksum over selected header fields, such as an IPv4 checksum.
\end{itemize}

Switch implementations are expected to restrict action processing so that header modifications remain consistent with the declared packet format.

\subsection{Control program and mTag processing sequence}

The control program specifies the flow from one table to the next using functions, conditionals, and table references. The example's logical sequence is:
\begin{enumerate}
  \item \textbf{Source check:} check consistency between the received packet and ingress port. An mTag should appear only on a port connected to a core switch. If an mTag is present, strip it and record that fact in metadata.
  \item \textbf{Local switching:} attempt to deliver the packet to a locally connected end host.
  \item \textbf{mTag table:} if local switching does not resolve an egress port, select/apply an mTag for core forwarding.
  \item \textbf{Egress check:} verify that the egress state is resolved and that the packet has not been incorrectly retagged. An unknown egress can trigger a notification to the SDN control stack.
\end{enumerate}

Metadata is important here: it records information such as whether the incoming packet carried an mTag and whether the ingress check found an error. Later tables can use this information to avoid retagging packets that should not be retagged.

\section{Compiling a P4 program}

The compiler maps the target-independent P4 description to the target switch's hardware or software, allocates resources, and generates device configuration. The compilation strategy differs according to the target's capabilities.

\subsection{Compiling packet parsers}

For a switch with a programmable parser, the compiler translates the parser description into a parsing state machine. For a fixed-parser switch, the compiler checks whether the requested parser is consistent with the parser the target already implements. A parser state-table entry specifies the current state, the field value being matched, and the next state. For example, in the mTag parser, VLAN EtherType \texttt{0xaaaa} transitions to the mTag state, while EtherType \texttt{0x800} transitions to IPv4; other values can lead to a stop state.

\subsection{Compiling control programs and table dependencies}

An imperative control-flow program is convenient for describing the intended logical behavior, but does not explicitly expose all table dependencies or parallelization opportunities. The compiler analyzes the program, identifies dependencies, constructs a TDG, and then maps that graph to the target's resources. This is the same two-step process introduced earlier: P4 program $\rightarrow$ TDG and dependency analysis $\rightarrow$ target-specific configuration.

The paper discusses several target cases:
\begin{itemize}
  \item \textbf{Software switches:} parsing, number of tables, and table configuration are controlled by software, so the compiler can map the mTag table graph directly to software tables. Table type information constrains table width, height, and matching criteria (exact, prefix, or wildcard); software data structures can optimize ternary or prefix matching.
  \item \textbf{Hardware switches with RAM and TCAM:} RAM and hashing can implement efficient exact matching, such as the edge mTag table. A core table that matches only a subset of tag bits can be mapped to TCAM.
  \item \textbf{Switches with parallel tables:} dependency analysis can place independent tables in parallel and dependent tables in series. In the example, local switching and the mTag table may execute in parallel until the action that sets the mTag is reached.
  \item \textbf{Switches that apply actions at the end of the pipeline:} intermediate stages can record intended modifications in metadata, and the final stage performs the actual writes. Metadata can indicate whether the mTag should be added or removed.
  \item \textbf{Switches with few physical tables:} the compiler can combine multiple logical P4 tables into fewer physical tables. For example, local switching and mTag lookup may be combined. When rules are installed at runtime, a rule translator can compose rules from the two logical tables into entries for one physical table.
\end{itemize}

These cases illustrate target independence: the P4 program states the desired behavior, while the compiler adapts the implementation to the resources and execution model of the chosen switch.

\section{Conclusion and key distinctions}

The paper presents P4 as a first-step proposal for more flexible switches and as a possible direction for an ``OpenFlow 2.0''-style interface. Its central argument is that a control plane should be able to describe how packets are parsed and processed, rather than being limited to a fixed set of header fields and predefined actions. A compiler can translate that description into a TDG, analyze dependencies, and map the result onto different switch targets, including optimized hardware.

The authors explicitly describe the work as a strawman proposal intended to contribute to discussion. Several aspects remain undefined in the proposal, including congestion-control primitives, queueing disciplines, and traffic monitoring. The proposed approach is a configuration language plus compilers that generate low-level target-specific configurations.

\subsection{Revision checklist}

\begin{itemize}
  \item \textbf{P4 is not simply a replacement for OpenFlow.} P4 describes/configures packet parsing and processing; an SDN control protocol such as OpenFlow can still populate the resulting forwarding tables.
  \item \textbf{Configure versus populate:} configuration defines the parser and table structure; population installs/removes entries that determine current policy.
  \item \textbf{Protocol independence versus target independence:} protocol independence allows new packet formats and named fields; target independence lets the same logical program be compiled for different hardware/software targets.
  \item \textbf{Parser versus match+action tables:} the parser extracts fields and determines the recognized header sequence; tables match those fields and execute actions.
  \item \textbf{Ingress versus egress:} ingress chooses forwarding/replication/drop/queue behavior; egress performs per-instance packet modifications.
  \item \textbf{P4 control flow versus TDG:} the programmer writes imperative control logic; the compiler derives a graph of dependencies to expose safe parallelism and guide placement.
  \item \textbf{mTag motivation:} hierarchical four-byte label fields let core switches forward using a relevant tag byte, reducing dependence on large Layer 2 destination tables.
  \item \textbf{Why a compiler matters:} different targets have different parsers, memories, pipeline structures, table counts, and action timing; the compiler adapts one target-independent description to those constraints.
\end{itemize}